% =====================================================================
%  PLANTILLA — copia este archivo como equipo_NN.tex por cada equipo
%  observado (todos menos el tuyo) y agrégalo en informe.tex con \input.
%
%  Reglas de la guía para llenar:
%  - La evidencia debe ser OBSERVABLE: n.º de diapositiva, minuto,
%    concepto explicado u omitido, error técnico, salida de la demo,
%    pregunta y respuesta. No vale "estuvo bien".
%  - Nivel 4 solo si se cumple TODO. Si dudas entre dos niveles, el inferior.
%  - Toda evidencia sale de TU hoja de campo de ese equipo (foto en hojas_campo/).
%  - 3 o más dimensiones sin evidencia = CERO.
% =====================================================================

\section{\completar{Equipo N --- Tema}}

\identificacion
  {\completar{nombre o número del equipo}}
  {\completar{T? --- nombre del tema}}
  {\completar{dd/mm/2026}}
  {\completar{hh:mm}}
  {\completar{hh:mm}}

\begin{dimensiones}
  % D1 Cumplimiento de la guía (10): conceptos clave, demo en vivo,
  %    comparación y recomendación.
  \dimension{D1 Cumplimiento de la guía (10)}
    {\faltanivel}
    {\completar{qué partes de la guía cubrieron y cuáles faltaron}}

  % D2 Fundamentación conceptual (20): nivel 4 = al menos 5 conceptos clave
  %    sin errores y acorde a la fuente oficial (RFC, OASIS, W3C, Jakarta EE...).
  %    Errores graves que bajan D2: WebSocket sobre UDP, JMS como protocolo,
  %    gRPC con JSON, TCP que conserva límites de mensaje,
  %    concurrencia = paralelismo, GET para modificar datos.
  \dimension{D2 Fundamentación conceptual (20)}
    {\faltanivel}
    {\completar{conceptos explicados (diapositiva/minuto) y errores detectados}}

  % D3 Desarrollo técnico (15): nivel 4 = demo en vivo en procesos distintos
  %    con un caso de fallo + comparación de >= 2 alternativas con >= 3
  %    criterios medibles.
  \dimension{D3 Desarrollo técnico (15)}
    {\faltanivel}
    {\completar{qué mostró la demo, si hubo caso de fallo, qué comparó}}

  % D4 Análisis y recomendación (20): nivel 4 = cuándo usar y cuándo no,
  %    >= 2 razones ligadas a un caso concreto, reconoce limitaciones.
  \dimension{D4 Análisis y recomendación (20)}
    {\faltanivel}
    {\completar{recomendación dada, razones, caso, limitaciones}}

  % D5 Comunicación (25): no lee, tiempo T +- 2 min, responde con precisión
  %    una pregunta técnica, diapositivas legibles, exponen todos.
  \dimension{D5 Comunicación (25)}
    {\faltanivel}
    {\completar{duración real, pregunta hecha y respuesta, quiénes expusieron}}

  % D6 Fuentes (10): fuente visible en cada dato/figura, fuentes primarias
  %    o indexadas, formato IEEE.
  \dimension{D6 Fuentes (10)}
    {\faltanivel}
    {\completar{diapositivas con/sin fuente, tipo de fuentes, formato}}
\end{dimensiones}

% Pon aquí los mismos niveles de la tabla, en orden D1..D6 (solo números):
\notasugerida{0}{0}{0}{0}{0}{0}

\retroalimentacion
  {\completar{una fortaleza concreta, con la evidencia que la respalda}}
  {\completar{una mejora accionable: qué hacer distinto y cómo}}

\clearpage
